\documentclass[10pt,twocolumn,letterpaper]{article}

% --- Geometry ---
\usepackage[
  top=0.75in, bottom=1in, left=0.75in, right=0.75in,
  columnsep=0.3in
]{geometry}

% --- Fonts and encoding ---
\usepackage[T1]{fontenc}
\usepackage{mathptmx}          % Times Roman for body
\usepackage{courier}           % Courier for monospace
\usepackage[scaled=0.92]{helvet} % Helvetica for sans-serif

% --- Math and symbols ---
\usepackage{amsmath,amssymb}

% --- Tables and figures ---
\usepackage{booktabs}
\usepackage{tabularx}
\usepackage{array}

% --- Bibliography ---
\usepackage[numbers,sort&compress]{natbib}

% --- Hyperlinks ---
\usepackage[
  colorlinks=true,
  linkcolor=black,
  citecolor=black,
  urlcolor=blue!70!black,
  bookmarks=false
]{hyperref}

% --- Misc ---
\usepackage{enumitem}
\usepackage{xcolor}
\usepackage{microtype}
\usepackage{balance}           % Balance last page columns

% --- Section formatting ---
\usepackage{titlesec}
\titleformat{\section}{\normalfont\large\bfseries}{\thesection.}{0.5em}{}
\titleformat{\subsection}{\normalfont\normalsize\bfseries}{\thesubsection}{0.5em}{}
\titlespacing*{\section}{0pt}{10pt plus 2pt minus 2pt}{4pt plus 1pt}
\titlespacing*{\subsection}{0pt}{8pt plus 2pt minus 1pt}{2pt plus 1pt}

% --- Title formatting ---
\makeatletter
\renewcommand{\maketitle}{%
  \twocolumn[%
    \begin{center}
      {\LARGE\bfseries\@title\par}
      \vskip 1em
      {\large\@author\par}
      \vskip 0.5em
      {\normalsize\itshape\@date\par}
      \vskip 1.5em
    \end{center}
  ]%
}
\makeatother

% --- Abstract formatting ---
\renewenvironment{abstract}{%
  \small
  \begin{center}
    \textbf{Abstract}\vspace{-0.5em}
  \end{center}
  \begin{quotation}\noindent
}{%
  \end{quotation}\vspace{0.5em}\par
}

% --- Compact lists ---
\setlist{nosep,leftmargin=1.2em}

% --- Footnotes ---
\renewcommand{\footnotesize}{\scriptsize}

% --- Code formatting ---
\newcommand{\code}[1]{\texttt{\small #1}}

\begin{document}

\title{Eigenius: A Typed Knowledge Graph with Epistemic Status\\
  Tracking for AI-Driven Science and Engineering}

\author{%
  Hans-Martin Will\\
  \texttt{hwill@acm.org}
}

\date{}

\maketitle
\thispagestyle{empty}

% ============================================================
\begin{abstract}
Large language models produce text that reads like knowledge but carries no epistemic warranty.
A correct derivation and a fluent hallucination are structurally indistinguishable in the output.
We present Eigenius, an open-source platform that addresses this problem by unifying structured knowledge representation, typed processing pipelines, LLM integration, and reasoning trace capture within a single self-describing knowledge graph.
Every resource in the system---domain facts, ontology definitions, program specifications, LLM interaction records, and formal proofs---is a typed resource with tracked provenance and queryable epistemic status.
A dependent type theory (Mini-TT) provides static guarantees for processing pipelines: termination, type safety, and exhaustive error handling.
Four epistemic categories---\emph{declared}, \emph{observed}, \emph{derived}, and \emph{verified}---are computed from provenance chains and enforced by the type system, making the boundary between verified knowledge and plausible-sounding output a structural property of the system rather than a matter of trust.
We describe the architecture, report on an implementation covering the kernel, query language, type checker, and persistent storage, and discuss the design's implications for neurosymbolic systems that must maintain epistemic integrity at scale.
\end{abstract}


% ============================================================
\section{Introduction}
\label{sec:intro}

The integration of large language models (LLMs) into scientific and engineering workflows has created a new category of epistemic risk.
An LLM can produce a convincing paragraph about quantum error correction or protein folding dynamics that is subtly wrong, and nothing in the output's structure reveals the error.
The confidence sounds the same whether the reasoning is correct or confabulated.
For frontier research---drug discovery, quantum physics, autonomous systems, complex infrastructure---this ambiguity is not an acceptable trade-off.

Existing infrastructure routes around this problem rather than solving it.
Orchestration frameworks (LangChain~\cite{langchain}, Semantic Kernel~\cite{semantickernel}) treat LLM outputs as opaque strings with no type discipline governing inter-step data flow.
Knowledge graph platforms (Neo4j, Neptune) store conclusions but not the derivations that produced them.
RAG systems ground retrieval but not reasoning.
Formal verification systems (Lean~\cite{lean4}, Coq) provide the strongest guarantees but exist in a separate universe from LLM-driven workflows.

We present \textbf{Eigenius}, a platform whose core architectural commitment is that every piece of knowledge---every fact, derivation, LLM interaction, and formal proof---is represented as a \emph{typed resource} in a unified knowledge graph with tracked provenance and queryable epistemic status.
The system maintains four epistemic categories:
\emph{declared} (human assertions),
\emph{observed} (recorded facts with provenance),
\emph{derived} (computed by typed pipelines with full audit trails),
and \emph{verified} (carrying machine-checked proofs).
These are not labels---they are computed from provenance chains and enforced structurally.

The key technical contributions are:
\begin{enumerate}
  \item A self-describing data model (\emph{Eigon}) where classes, properties, programs, and reasoning traces are all uniform typed resources;
  \item A query language (\emph{EigenQL}) based on stratified Datalog with aggregation that operates uniformly over domain data and meta-level reasoning traces;
  \item A program model grounded in Mini-TT~\cite{minitt} dependent type theory providing static guarantees for processing pipelines;
  \item A reasoning trace schema where traces mirror the expression tree, serving simultaneously as memoization cache, provenance record, and epistemic status computation substrate.
\end{enumerate}


% ============================================================
\section{The Eigon Data Model}
\label{sec:eigon}

Eigon is the canonical typed data format, derived from Atomic Data~\cite{atomicdata} and adapted to the requirements of a knowledge graph with epistemic tracking.
Atomic Data's core ideas---typed properties identified by URLs, self-describing schemas, and a JSON serialization where every value has a known type---provide the foundation.
Eigon retains this philosophy but makes several structural changes:
(i) identifiers are IRIs using the \code{urn:} scheme (RFC~3987) rather than fetchable HTTP URLs, because type resolution comes from the loaded ontology in the layer stack, not from dereferencing the identifier over the network;
(ii) there is no \code{@context} mechanism---namespace resolution is handled by the layer system;
(iii) class membership is expressed via an \code{is\_a} array (supporting multiple class membership), rather than Atomic Data's single-valued \code{is-a};
(iv) property ownership is inverted---classes declare which properties they require, rather than properties declaring which class they belong to;
and (v) a three-layer type system (primitive data types, format constraints, and content types) replaces Atomic Data's single-level typing.

All data---ontology definitions, instance resources, processing pipelines, reasoning traces---is represented as \emph{resources}.
There are no separate internal representations for different resource kinds; everything is a uniform \code{Resource} value distinguished by its \code{is\_a} property.
This uniformity is the mechanism that makes the system coherently queryable across meta and object levels.

\subsection{Resource Identity and Structure}

A top-level resource has a globally unique identity expressed as an IRI in the \code{@id} field, the only reserved key in Eigon-JSON.
All property keys are full IRIs---there are no abbreviated forms in the core format.
Class membership is expressed via the \code{is\_a} property, always an array supporting multiple class membership.
Embedded resources (nested objects without \code{@id}) are addressable only through their parent.

The type system separates three concerns on property values, inspired by JSON Schema:
\emph{primitive data types} determine JSON-level representation (string, integer, float, boolean, resource, resource\_array, value\_array, json);
\emph{formats} constrain string values (date, datetime, IRI, UUID, regex);
and \emph{content types} declare the MIME type of embedded content (text/markdown, image/png).
This separation allows precise typing without an unbounded proliferation of primitive types.

\subsection{Self-Description}

The Core Ontology is self-describing: \code{Class} is an instance of \code{Class}, \code{Property} is an instance of \code{Class}, \code{is\_a} is an instance of \code{Property}.
This circularity is resolved by hardcoding the Core Ontology in the kernel.
Once loaded, all subsequent ontology processing uses the standard validation pipeline.

Self-description means the system can query its own structure with the same language used for domain data.
``What classes exist?'' is the same kind of query as ``what molecules bind to this receptor?''
The meta-level and object-level share a single query language, type system, and storage layer.

\subsection{Decidability Boundary}

The Core Ontology is deliberately bounded to constructs where type-checking and query evaluation remain decidable given a finite ontology.
Transitive properties, union/intersection class constructors, negation/complement, and property chains are excluded.
The guiding principle: the core ontology \emph{validates}; registered capabilities \emph{reason}.
Any construct requiring the kernel to derive new facts belongs in the capability layer.


% ============================================================
\section{Immutable Layer System}
\label{sec:layers}

The ontology within any execution context is organized as a linear stack of immutable layers.
Each layer has exactly one parent.
The Core Ontology is the fixed root.

Resource resolution is a linear scan from the top of the stack to the root, stopping at the first definition found.
Committed layers are immutable---once a layer becomes a parent, it cannot be modified.
This provides snapshot isolation structurally: a snapshot is a pointer to a top layer, and the entire history below is immutable.
Layer identifiers are content-addressed (SHA-256 of CBOR deterministic encoding~\cite{rfc8949}).

The global structure across all execution contexts is a tree---analogous to Git branches.
Each context sees a linear history, but contexts may diverge above a shared ancestor.
Conflict detection uses optimistic concurrency: first commit wins.

Namespace governance is enforced at the layer level.
Each layer declares a namespace prefix; resources must have IRIs within that namespace.
Cross-namespace writes are structurally impossible, while cross-namespace reads (references) are unrestricted.
The \code{urn:eigenius:core:} namespace is protected---attempts to write resources there from any non-core layer are rejected at ingestion time.


% ============================================================
\section{EigenQL: Typed Semantic Queries}
\label{sec:eigenql}

EigenQL is a typed stratified Datalog with aggregation.
It supports conjunctive queries, recursive rule definitions (\code{DEFINE}), stratified negation, \code{GROUP BY} with aggregation (COUNT, SUM, AVG, MIN, MAX), \code{ORDER BY}, \code{LIMIT}/\code{OFFSET}, and \code{DISTINCT}.

A program consists of zero or more \code{DEFINE} clauses followed by a query.
Typed \code{MATCH} patterns constrain matches to class instances (including subclasses); properties referenced in typed patterns are validated against the ontology schema.
Variables are strongly typed based on their binding context, and type checking is performed at query submission before evaluation begins.

Recursive rules use bottom-up seminaive fixpoint evaluation~\cite{alice}.
Negated patterns are subject to stratification checking---the system builds a dependency graph and rejects programs with negation cycles before evaluation.
Queries without negation are monotonic (adding resources can only increase results); queries with negation are flagged for full re-evaluation on layer changes.

Because programs, traces, and domain data are all typed resources, EigenQL can query across them uniformly.
For example: ``Which LLM calls consumed more than 500 tokens and used the Anthropic provider?'' is a standard EigenQL query over \code{ComponentTrace} resources in the same knowledge graph as domain data.


% ============================================================
\section{Typed Program Model}
\label{sec:programs}

Programs in Eigenius are typed expressions in a functional language grounded in Mini-TT~\cite{minitt}, a minimal dependent type theory that is a direct fragment of CIC---the Calculus of Inductive Constructions underlying Lean~4~\cite{lean4}.
A program that passes type checking carries formal guarantees: it terminates, is type-safe, and produces output of the declared type.

\subsection{Expression Forms}

Programs are represented as Eigon resources---the same format as everything else.
Twelve expression forms (Program, Let, Apply, Var, Lambda, Case, Pair, Construct, Project, Map, Reduce, Literal) map directly to Mini-TT terms with no translation layer.
Components are typed functions registered in the ontology; calling one is function application.
Map and Reduce are language primitives (not library functions) so the scheduler has explicit knowledge of parallelism opportunities.

\subsection{Type Checking via NbE}

Type checking uses Normalization by Evaluation (NbE)~\cite{nbe} with bidirectional type checking.
Ontology classes are represented as dependent record types (nested $\Sigma$-types over required properties).
Type equality is decided by normalizing both sides and comparing normal forms.
Subclass compatibility is modeled as coercion functions computed from the ontology's subclass hierarchy at validation time.

Partial evaluation is a direct consequence of NbE, requiring no additional machinery.
When some inputs are bound concretely and others are abstract (generators), the evaluator reduces what it can and produces neutral terms for stuck computations.
The result is a well-typed residual program.
This enables pipeline specialization, static configuration resolution, and incremental revalidation when layers change.

\subsection{Components and Capabilities}

Components declare three capability levels: \emph{pure} (data transformation only), \emph{read} (layer chain access), and \emph{io} (network/LLM calls).
They also declare determinism and fallibility.
Fallible components return a labeled sum type $\mathsf{Sum}(\mathsf{ok} : A \mid \mathsf{err} : E)$; the type checker enforces exhaustive handling of both branches.

\begin{table}[t]
  \centering
  \small
  \caption{Component capability levels and their dispatch behavior.}
  \label{tab:capabilities}
  \begin{tabular}{@{}llll@{}}
    \toprule
    \textbf{Level} & \textbf{Executor} & \textbf{Trace check} & \textbf{Parallel} \\
    \midrule
    Pure & Kernel (local) & No & Yes \\
    Read & Kernel (local) & No & Yes \\
    IO   & Orchestrator   & Yes (cached) & Yes \\
    \bottomrule
  \end{tabular}
\end{table}


% ============================================================
\section{Execution and Reasoning Traces}
\label{sec:traces}

The Rust kernel executes programs by walking the expression tree.
Pure expressions evaluate locally; IO component applications dispatch to a Deno/TypeScript orchestrator.
The key architectural insight is that \emph{reasoning traces serve simultaneously as memoization cache, provenance record, and epistemic status substrate}.

\subsection{Trace-Based Durability}

Before dispatching any IO component, the kernel checks for a valid trace with matching inputs.
If found, the cached output is used without re-execution.
Each completed component call produces a trace stored as a typed Eigon resource recording the component, input/output hashes, provider, model, token usage, and latency.

This provides durability (crash recovery resumes from the last completed trace), caching (re-execution with identical inputs is instant), and incremental re-evaluation (changing one component's prompt re-executes only that step).
The connection to NbE is formal: a trace is a \emph{known value}, a missing trace is a \emph{neutral term} (blocked on an unknown).
Trace-based execution is partial evaluation where ``known'' means ``previously computed and traced.''

\subsection{Trace Structure Mirrors Expression Structure}

The trace type system mirrors the expression language.
Each expression form has a corresponding trace form (LetTrace, ComponentTrace, MapTrace, CaseTrace, etc.).
Each evaluation step returns both its result and its trace; the trace tree builds up naturally as evaluation recurses---no separate tracing pass is needed.

Only \code{ComponentTrace}s (IO computation with external effects) receive content-addressed identifiers for memoization; other trace types are embedded in the tree.

\subsection{Epistemic Status from Provenance}

Four epistemic categories are computed from the resource's provenance chain:

\begin{description}[style=unboxed,leftmargin=0em]
  \item[Declared] --- human-authored axioms, definitions, prompt templates. Anchored in intent, not computation.
  \item[Observed] --- facts recorded from external reality with tracked provenance. The system vouches for provenance, not truth.
  \item[Derived] --- computed by typed programs with full audit trails. The trace tree links every output to every input.
  \item[Verified] --- carries a machine-checked formal proof. The proof term is itself a typed resource.
\end{description}

\noindent These are enforced via base classes (\code{DeclaredResource}, \code{ObservedResource}, \code{DerivedResource}, \code{VerifiedResource}) with required provenance properties.
Transitions are monotonic upward---a derived result can be promoted to verified by attaching a proof, but nothing can be silently downgraded.

Universe stratification prevents self-referential paradox: a trace at level $N$ can only reference resources at level $N{-}1$ or below.
This is enforced at ingestion time.
Queries may join across levels---stratification governs \emph{authorship}, not \emph{visibility}.

\begin{table}[t]
  \centering
  \small
  \caption{Epistemic categories, their trace types, and enforcement.}
  \label{tab:epistemic}
  \begin{tabular}{@{}lll@{}}
    \toprule
    \textbf{Status} & \textbf{Trace type} & \textbf{Required property} \\
    \midrule
    Declared & DeclarationTrace & \code{declared\_by} \\
    Observed & ObservationTrace & \code{source} \\
    Derived  & ProgramTrace     & \code{derivation} \\
    Verified & VerificationTrace & \code{verification} \\
    \bottomrule
  \end{tabular}
\end{table}


% ============================================================
\section{Implementation Status}
\label{sec:impl}

The system is implemented as a Rust kernel with a Deno/TypeScript orchestration layer.
Four of six planned phases are complete:

\textbf{Phase~0 (Foundation).}
The uniform resource model, Core Ontology (self-describing bootstrap), immutable layer system with content-addressed identifiers, execution contexts, in-memory storage backend, namespace enforcement, and a CLI with a \code{load} command.

\textbf{Phase~1 (Query).}
Full EigenQL implementation: lexer, parser, type checker, and evaluator supporting \code{DEFINE} with recursive rules, seminaive fixpoint evaluation, stratified negation, aggregation, \code{GROUP BY}, \code{ORDER BY}, \code{LIMIT}/\code{OFFSET}, \code{DISTINCT}, and dot-path navigation.

\textbf{Phase~2 (Pipelines).}
Mini-TT dependent type theory ported from the Haskell reference implementation.
NbE evaluator with bidirectional type checking.
Program ontology comprising 54 resources.
End-to-end pipeline: parse, type-check, execute.

\textbf{Phase~3 (Service).}
gRPC server (tonic), CBOR serialization for storage and wire, RocksDB persistent storage, CLI dual mode (embedded kernel or remote gRPC client), server integration tests.

Remaining phases cover LLM integration with reasoning trace recording (Phase~4), a human-friendly surface syntax (Phase~4.5), and WASM-sandboxed extensibility for untrusted domain capabilities (Phase~5).
A parallel Lean~4 formal track specifies the metatheory.


% ============================================================
\section{Related Work}
\label{sec:related}

\textbf{LLM orchestration.}
LangChain~\cite{langchain}, LlamaIndex~\cite{llamaindex}, and Semantic Kernel~\cite{semantickernel} compose LLM calls into workflows but treat outputs as opaque strings.
There is no type system governing inter-step data flow and no mechanism to distinguish verified from unverified conclusions.
Eigenius provides both through its dependent type system and epistemic categories.

\textbf{Knowledge graphs.}
RDF~\cite{rdf}, OWL~\cite{owl}, and platforms like Neo4j and Stardog provide structured storage and querying.
However, they are disconnected from reasoning processes: a graph can store a conclusion but not the replayable, auditable derivation that produced it.
Eigenius's reasoning traces close this gap.

\textbf{Atomic Data.}
The Eigon data model is derived from Atomic Data~\cite{atomicdata}, which pioneered the approach of typed, self-describing resources with JSON serialization where every property value has a known type.
The key adaptations for Eigenius are described in Section~\ref{sec:eigon}: non-fetchable IRIs in place of URLs, class-declared property ownership, multiple class membership, and a three-layer type system.
These changes reflect the requirements of an epistemic knowledge graph operating over an immutable layer stack rather than a linked-data web of individually fetchable resources.

\textbf{Provenance systems.}
The W3C PROV model~\cite{prov} provides a vocabulary for provenance.
Eigenius's approach is more tightly integrated: traces are structural components of the knowledge graph itself, typed by the same ontology and queryable with the same language.

\textbf{Typed workflow systems.}
CWL~\cite{cwl} and WDL provide typed workflow descriptions but target computational biology pipelines, not general AI-driven reasoning with epistemic tracking.
Eigenius's dependent types provide stronger guarantees (termination, partial evaluability) than the simple type systems in workflow languages.

\textbf{Dependent types in practice.}
Lean~4~\cite{lean4}, Agda~\cite{agda}, and Idris~\cite{idris} use dependent type theories for formal verification and programming.
Eigenius uses Mini-TT~\cite{minitt}---a minimal fragment of CIC---not as a general programming language but as the foundation for validating AI processing pipelines against ontology schemas.

\textbf{Neurosymbolic architectures.}
DeepProbLog~\cite{deepproblog}, NeurASP~\cite{neurasp}, and Logic Tensor Networks~\cite{ltn} integrate neural and symbolic computation, typically by embedding differentiable computation within logic programs.
Eigenius takes a different approach: LLMs are black-box components in typed pipelines, and the symbolic layer provides structure, provenance, and epistemic discipline around their outputs rather than integrating at the gradient level.


% ============================================================
\section{Discussion}
\label{sec:discussion}

\subsection{Verification as a Spectrum}

The four epistemic categories represent a deliberate spectrum.
The system is useful before any formal proofs exist---recording observations and running typed pipelines with reasoning traces is valuable on its own.
Formal verification deepens incrementally.
The architecture does not demand that everything be proved; it demands that the system always show what \emph{has} been proved and what has not.

\subsection{The Self-Describing Property}

Self-description is not an academic curiosity.
It means the system can query its own programs, traces, and ontology with the same language used for domain data.
This enables queries like ``which processing programs reference a specific LLM provider?'' or ``what unverified assumptions does this conclusion depend on?'' without special-purpose tooling.

\subsection{Limitations}

The dependent type system guarantees structural properties of programs (termination, type safety) but not semantic properties of LLM outputs.
A well-typed pipeline that invokes an LLM is guaranteed to produce output of the declared type, but the content of that output reflects the LLM's statistical patterns, not logical validity.
Eigenius makes this distinction explicit rather than resolving it---the gap between \emph{derived} and \emph{verified} is a permanent feature of the architecture.

The system's overhead---mandatory provenance, typed traces, CBOR serialization---adds latency and storage cost.
For applications where epistemic tracking is unnecessary, lighter-weight orchestration frameworks are more appropriate.


% ============================================================
\section{Conclusion}
\label{sec:conclusion}

We have presented Eigenius, a platform that unifies typed knowledge representation, dependent-type-checked processing pipelines, LLM integration, and reasoning trace capture within a single self-describing knowledge graph.
Its core contribution is making the epistemic status of AI-assisted conclusions---declared, observed, derived, or verified---a structural, queryable, and enforceable property of the system.
Four implementation phases are complete, covering the data model, query language, type checker, and persistent storage service.
The architecture provides a concrete path from informal LLM-assisted analysis to formally verified conclusions, with the system always showing where the boundary lies.

\balance

% ============================================================
\bibliographystyle{plainnat}
\begin{thebibliography}{16}

\bibitem[Coquand et~al.(2009)]{minitt}
T.~Coquand, Y.~Kinoshita, B.~Nordstr{\"o}m, and M.~Takeyama.
\newblock A simple type-theoretic language: {Mini-TT}.
\newblock In \emph{From Semantics to Computer Science: Essays in Honour of Gilles Kahn}, pages 139--164. Cambridge University Press, 2009.

\bibitem[de~Moura and Ullrich(2021)]{lean4}
L.~de~Moura and S.~Ullrich.
\newblock The {Lean~4} theorem prover and programming language.
\newblock In \emph{CADE-28}, LNCS 12699, pages 625--635. Springer, 2021.

\bibitem[Abel et~al.(2007)]{nbe}
A.~Abel, T.~Coquand, and P.~Dybjer.
\newblock Normalization by evaluation for {Martin-L{\"o}f} type theory with typed equality judgements.
\newblock In \emph{LICS 2007}, pages 3--12. IEEE, 2007.

\bibitem[Chase(2023)]{langchain}
H.~Chase.
\newblock {LangChain}.
\newblock \url{https://langchain.com/}, 2023.

\bibitem[Liu(2023)]{llamaindex}
J.~Liu.
\newblock {LlamaIndex}.
\newblock \url{https://llamaindex.ai/}, 2023.

\bibitem[Microsoft(2023)]{semantickernel}
Microsoft.
\newblock {Semantic Kernel}.
\newblock \url{https://github.com/microsoft/semantic-kernel}, 2023.

\bibitem[Cyganiak et~al.(2014)]{rdf}
R.~Cyganiak, D.~Wood, and M.~Lanthaler.
\newblock {RDF 1.1} concepts and abstract syntax.
\newblock W3C Recommendation, 2014.

\bibitem[Hitzler et~al.(2012)]{owl}
P.~Hitzler, M.~Kr{\"o}tzsch, B.~Parsia, P.F.~Patel-Schneider, and S.~Rudolph.
\newblock {OWL~2} web ontology language primer.
\newblock W3C Recommendation, 2012.

\bibitem[Meindertsma(2020)]{atomicdata}
J.~Meindertsma.
\newblock {Atomic Data: A modular specification for sharing, modifying and modeling graph data}.
\newblock W3C Community Group Specification, 2020--2026.
\newblock \url{https://docs.atomicdata.dev/}.

\bibitem[Groth and Moreau(2013)]{prov}
P.~Groth and L.~Moreau.
\newblock {PROV-Overview}.
\newblock W3C Working Group Note, 2013.

\bibitem[Amstutz et~al.(2022)]{cwl}
P.~Amstutz et~al.
\newblock Common workflow language, v1.2.
\newblock \url{https://www.commonwl.org/}, 2022.

\bibitem[Norell(2009)]{agda}
U.~Norell.
\newblock Dependently typed programming in {Agda}.
\newblock In \emph{AFP 2008}, LNCS 5832, pages 230--266. Springer, 2009.

\bibitem[Brady(2021)]{idris}
E.~Brady.
\newblock {Idris~2}: Quantitative type theory in practice.
\newblock In \emph{ECOOP 2021}, LIPIcs 194, 2021.

\bibitem[Manhaeve et~al.(2018)]{deepproblog}
R.~Manhaeve, S.~Dumancic, A.~Kimmig, T.~Demeester, and L.~De~Raedt.
\newblock {DeepProbLog}: Neural probabilistic logic programming.
\newblock In \emph{NeurIPS 2018}, 2018.

\bibitem[Yang et~al.(2020)]{neurasp}
Z.~Yang, A.~Ishay, and J.~Lee.
\newblock {NeurASP}: Embracing neural networks into answer set programming.
\newblock In \emph{IJCAI 2020}, 2020.

\bibitem[Badreddine et~al.(2022)]{ltn}
S.~Badreddine, A.~Garcez, L.~Serafini, and M.~Spranger.
\newblock Logic tensor networks.
\newblock \emph{Artificial Intelligence}, 303, 2022.

\bibitem[Abiteboul et~al.(1995)]{alice}
S.~Abiteboul, R.~Hull, and V.~Vianu.
\newblock \emph{Foundations of Databases}.
\newblock Addison-Wesley, 1995.

\bibitem[Bormann and Hoffman(2020)]{rfc8949}
C.~Bormann and P.~Hoffman.
\newblock Concise binary object representation ({CBOR}).
\newblock RFC 8949, IETF, 2020.

\end{thebibliography}

\end{document}